\section{A uniform maximal-slope bound for the principal Lie bundle}
\label{sec:sr-general-principal-slope}

All degrees in this section are divided by the
degree of the number field.
For a Hermitian bundle \(\overline V\), write
\(\widehat{\deg}_{\mathrm n}\overline V\) for this
normalized degree and
\(\widehat\mu(\overline V)=
\widehat{\deg}_{\mathrm n}\overline V/\operatorname{rk}V\).
The metric on the Lie bundle is the metric
\(\mathrm{HodgeIP}\) of the principal polarization.
In period coordinates its squared norm is
\(\pi H(v,v)\), where \(H\) is the positive
Hermitian Riemann form.

\begin{srtheorem}\label{thm:sr-general-principal-slope}
There is an absolute real constant \(C\ge0\)
with the following property.
Let \(L\) be any number field.
Let \(\mathcal G/\mathcal O_L\) be a smooth,
separated, commutative group scheme of finite
type with geometrically connected semiabelian
fibres, whose generic fibre \(A/L\) is an
abelian variety.
Suppose a symmetric, rigidified cubical line
bundle \(\mathcal L\) on \(\mathcal G\) induces
a principal ample polarization on \(A\).
Let
\[
 B=\operatorname{Lie}(\mathcal G/\mathcal O_L),
 \qquad g=\operatorname{rk}_{\mathcal O_L}B,
 \qquad
 h=\max\{-\widehat{\deg}_{\mathrm n}\overline B,0\}.
\]
Give \(B\) the preceding Hodge metric.
Then
\[
 \widehat\mu_{\max}(\overline B)
           \le C\log(2+g+h).                         \tag{SB1}
\]
The same \(C\) works with any compatible
endomorphism action and any CM host.
In particular the statement is exactly the
inequality required by \(\mathrm{SlopeBoundSt}\)
for every \(\mathrm{PPData}\) packet.
\end{srtheorem}

We prove the theorem below.
We retain the high directional-jet determinant
argument of the frozen source.
For its arithmetic theta input we use a
precise classical integral model theorem,
and check its compatibility with the given
connected model.
Thus the proof does not require an estimate
for the degree of a field used to acquire
torsion points or to principalize an ideal.

\subsection{The connected model and the normalization of the height}

\begin{srlemma}\label{lem:sr-slope-model-metric}
After any finite extension \(L'/L\), the
base change of \(\mathcal G\) is the identity
component of the Neron model of \(A_{L'}\).
Its rigidified line bundle and Lie metric
agree with the restrictions of any
Moret--Bailly model of the same rigidified
generic line bundle.
Moreover
\[
 h_{\mathrm F}^{\mathrm D}(A)
       =-\widehat{\deg}_{\mathrm n}\overline B,        \tag{SB2}
\]
where the superscript denotes Deligne's
normalization
\[
 \|\omega\|_{\mathrm D,\sigma}^{\,2}
   ={i^{g^2}\over(2\pi)^g}
                  \int_{A_\sigma(\mathbb C)}
                         \omega\wedge\overline\omega .
\]
All these identifications commute with
further finite extensions.
\end{srlemma}

\begin{proof}
The model assertion is
Bosch--Lutkebohmert--Raynaud,
\emph{Neron Models}, Section 7.4, Proposition 3
and Corollary 4.
These results apply over the localization at
every nonzero prime of \(\mathcal O_L\);
they identify a semiabelian model of an abelian
generic fibre with the identity component,
and allow faithfully flat extensions of
discrete valuation rings, including ramified ones.
They therefore identify the global connected
models as well.

Here is also the finite-place check on the
line bundle.
On the common connected smooth model,
the generic identification between two
rigidified extensions of the same line bundle
has a vertical Cartier divisor.
The model is regular, and every closed fibre
is geometrically integral and has multiplicity
one.
Consequently every vertical divisor is a sum
of whole fibres, hence is pulled back from
the base.
The rigidification at zero identifies this
base divisor with zero.
The resulting identification is integral.
The cubical and symmetric structures agree:
their discrepancies are units equal to one
on the dense generic fibre.
This proves the needed line-bundle
identification without making an assertion
about the extra components of a
Moret--Bailly model.

For (SB2), choose complex coordinates on
\(A_\sigma=\mathbb C^g/\Lambda_\sigma\).
If \(H_\sigma\) is the matrix of its Riemann
form, principality gives
\(\operatorname{vol}(\Lambda_\sigma)=
(\det H_\sigma)^{-1}\).
For a Lie basis with coordinate matrix \(W\),
the Gram determinant for \(B\) is
\[
             \det\operatorname{Gram}_{B,\sigma}
                =\pi^g|\det W|^2\det H_\sigma .
\]
The invariant differential dual to this basis
has Deligne norm squared
\[
 \pi^{-g}|\det W|^{-2}
                 \operatorname{vol}(\Lambda_\sigma)
             =(\det\operatorname{Gram}_{B,\sigma})^{-1}.
\]
The integral differential lattice is
\(\det(B)^\vee\).
Thus its Hermitian line is isometric to
\(\det(\overline B)^\vee\), proving (SB2).
The usual normalization without the factor
\(\pi^{-g}\) has height
\[
 h_{\mathrm F}^{0}(A)
       =h_{\mathrm F}^{\mathrm D}(A)
                         -{g\over2}\log\pi.           \tag{SB3}
\]
Normalized degrees and normalized sums over
complex embeddings are unchanged by the
finite extensions under consideration.
\end{proof}

A reference for the model assertions, with
the stated base-change scope, is
\href{https://www.math.stonybrook.edu/~kamenova/homepage_files/Bosch_Raynaud_Neron_Model_tc.pdf}
{Bosch--Lutkebohmert--Raynaud, Section 7.4,
Proposition 3 and Corollary 4}.
The complex normalization in (SB2) is the
one used in
\href{https://www.numdam.org/item/BSMF_2012__140_1_19_0.pdf}
{Pazuki, \emph{Theta height and Faltings height},
Section 2.2}.

\subsection{One integral theta section and its normalized translates}

We use the following exact part of the
Moret--Bailly model theorem.
For a symmetric ample \(L_A\) on an abelian
variety over \(\overline{\mathbb Q}\), and a
finite set \(F\) of algebraic points, there is,
after a finite extension of any given field
of definition, a semistable model carrying
a cubist Hermitian extension of \(L_A\) and
sections extending every point of \(F\).
For such models the integral \(L^2\) module of
global sections is independent of the model
under its generic identification and commutes
isometrically with finite base change.
When \(h^0(A,L_A)=1\), its normalized degree is
\[
 \widehat{\deg}_{\mathrm n}\pi_*\overline{\mathcal L}
       =-{1\over2}h_{\mathrm F}^{\mathrm D}(A)
                              -{g\over4}\log(2\pi).  \tag{SB4}
\]
These are Definition 3.1 and Theorem 3.4
(i), (iii), (iv), (v) of
\href{https://www.numdam.org/item/BSMF_2012__140_1_19_0.pdf}
{Pazuki, \emph{loc.\ cit.}}.
The model may have disconnected special
fibres. No good-reduction hypothesis is
imposed, and the metric uses Haar measure
of total mass one.

\begin{srlemma}\label{lem:sr-slope-joint-theta-witness}
After a finite extension of \(L\) there is an
integral nonzero theta section \(s\) on the
connected model with
\[
 b_\sigma=\|s\|_{2,\sigma}>0,\qquad
 \langle\log b_\sigma\rangle
       ={1\over2}h_{\mathrm F}^{\mathrm D}(A)
                               +{g\over4}\log(2\pi)
       \le {h\over2}+C_0g.                           \tag{SB5}
\]
Here and below brackets denote the normalized
sum over all complex embeddings.
For every positive integer \(m\), after a
further finite extension there are
\(r=m^{2g}\) integral sections \(s_P\) of
\(\mathcal L^{m^2}\), indexed by \(A[m]\),
which span its full complex space of sections
at every embedding and satisfy
\[
 \|s_P(z)\|_{\mathcal L^{m^2}}
                   =\|s(mz+P)\|_{\mathcal L}.        \tag{SB6}
\]
All comparisons with \(s\), \(B\), \(g\), and
the quantities in (SB5) are by base change
from the same original data.
\end{srlemma}

\begin{proof}
Choose a Moret--Bailly model and principalize
the rank-one projective module
\(\pi_*\mathcal L\) by a finite extension.
Let \(s\) be a generator.
Its restriction to the connected model is
integral by Lemma~\ref{lem:sr-slope-model-metric}.
Equation (SB4), and the degree formula for a
generator, give (SB5).
The section need not generate the larger
module of all sections on the connected
model. Neither the proof nor the jet
integrality below requires that assertion.

For a given \(m\), replace the model, after
finite extension, by one containing all
points of \(A[m]\) as sections.
Model independence and base-change
compatibility in the quoted theorem preserve
the integral section \(s\) under the fixed
generic identification.
We explain why its normalized translates
are integral; extending a torsion point
to the connected model is not being assumed.

For a torsion section \(P\), put
\(M_P=P^*\overline{\mathcal L}\), an integral
Hermitian line on the base.
The cube identity, including its isometric
rigidification, gives for every \(n\ge1\)
\[
 [n]^*\overline{\mathcal L}
   \simeq
 \overline{\mathcal L}^{\,n(n+1)/2}
       \otimes[-1]^*\overline{\mathcal L}^{\,n(n-1)/2}.
\]
Set \(A_m=m(m+1)/2\) and \(B_m=m(m-1)/2\).
Since \(mP=0\) as sections of the separated
group model, evaluation at \(P\) and \(-P\)
gives isometries
\[
 M_P^{A_m}\otimes M_{-P}^{B_m}\simeq\mathcal O,
 \qquad
 M_P^{B_m}\otimes M_{-P}^{A_m}\simeq\mathcal O.
\]
Eliminating \(M_{-P}\), and using
\(A_m^2-B_m^2=m^3\), gives
\[
                         M_P^{m^3}\simeq\mathcal O   \tag{SB7}
\]
isometrically and integrally.
This argument does not assume that a generic
symmetry of the line bundle extends over
every disconnected component.

A Hermitian integral line with an isometric
trivialization of a positive tensor power
becomes isometrically trivial after a finite
extension.
Indeed first principalize its underlying
ideal. If \(e\) is an integral generator
and \(e^k\) maps to a unit \(u\), adjoin a
\(k\)-th root of \(u\); the generator
\(e/u^{1/k}\) then has norm one at every
embedding. Apply this simultaneously to the
finitely many lines \(M_P\).

Let \(\mathcal G\) now denote the common
identity component.
The cube identity defines a biextension
\[
 \mathcal P(x,y)=
 \mathcal L(x+y)\otimes\mathcal L(x)^{-1}
                              \otimes\mathcal L(y)^{-1}.
\]
On \(\mathcal G\), the line in \(x\) obtained
by putting \(y=P\) satisfies
\([m]^*\mathcal P(x,P)=\mathcal P(x,mP)
\simeq\mathcal O\).
The integral identification of the symmetric
rigidified line on \(\mathcal G\) also gives
\([m]^*\mathcal L\simeq\mathcal L^{m^2}\).
Consequently
\[
 ([m]+P)^*\overline{\mathcal L}
       \simeq\overline{\mathcal L}^{m^2}\otimes M_P.
\]
These are isometries for the cubist metrics.
Pull back \(s\) and divide by the norm-one
integral generator of \(M_P\).
The result is \(s_P\); it is integral on
\(\mathcal G\), and the isometry proves (SB6).

For clarity, the spanning assertion is also
explicit.
In an Appell--Humbert trivialization write
\(s\) as a nonzero theta function \(\theta\).
If \(p\) represents \(P\in A[m]\), the
function representing \(s_P\) differs by a
nonzero scalar from
\[
 \theta(mz+p)\exp\{-\pi m H(z,p)\}.
\]
Its weighted norm is the norm on the right
of (SB6). Translation by \(\lambda/m\),
\(\lambda\in\Lambda\), acts, after removing
one common multiplier, by the character
\[
                \lambda\longmapsto
                         \exp(2\pi i\,\operatorname{Im}H(p,\lambda)).
\]
Different \(p\) modulo \(\Lambda\) give
different characters, by unimodularity
of the principal alternating form.
There are \(m^{2g}\) such functions, all
nonzero.
Eigenvectors for these different characters
are linearly independent.
Riemann--Roch gives
\(h^0(A,L_A^{m^2})=m^{2g}\), so they span.
Multiplication of the functions by the
nonzero normalization scalars does not
alter this conclusion.
\end{proof}

The same arithmetic construction also
explains the product-formula input behind
torsion sampling.
If \(s(P)\ne0\), evaluate it in the
trivialization (SB7).
The resulting scalar \(T_P\) is an algebraic
integer and
\[
             |\sigma T_P|=\|s(P)\|_\sigma^{m^3}.
\]
Hence
\(\langle\log\|s(P)\|_\sigma\rangle\ge0\).
When a complete packet of nonzero torsion
values is used, the common exponent \(m^3\)
cancels from the normalized sum.
The archived definition of \(\mathrm{TorsVal}\)
uses a more specific cubical trivialization
and the exponent \(m^2\).
We need neither identify that scalar with
\(T_P\) nor assert its universal integrality
to obtain (SB5)--(SB6).
This is a joint witness for the actual slope
argument, with all finite places included.

\subsection{The mean of the theta supremum}

\begin{srlemma}\label{lem:sr-slope-theta-mean-sup}
For the section of
Lemma~\ref{lem:sr-slope-joint-theta-witness}
there are positive \(M_\sigma\) such that
\[
 \|s\|_{\infty,\sigma}\le M_\sigma,\qquad
 \langle\log M_\sigma\rangle
        \le {h\over2}+C_1g\log(2+g+h).                \tag{SB8}
\]
The same bounds \(M_\sigma\), pulled back by
the embeddings, bound every section in (SB6),
for every \(m\).
Here \(C_1\) is absolute.
\end{srlemma}

\begin{proof}
Let \(\rho_\sigma\) be the first minimum
of the period lattice in the Riemann metric
\(H_\sigma\), without the factor \(\pi\).
The principal-polarization Gromov estimate
of Gaudron,
\href{https://lmbp.uca.fr/~gaudron/art19.pdf}
{\emph{Some explicit computations in Arakelov
geometry of abelian varieties}, Proposition 3.1},
is
\[
 \|s\|_{\infty,\sigma}
     \le b_\sigma(1.7g)^{g/2}
                       \max(1,\rho_\sigma^{-1})^{g/2}. \tag{SB9}
\]
This compares the same cubist norm and
the same probability-Haar \(L^2\) norm.

To verify its uniform averaged use, apply
Autissier's
\href{https://www.math.u-bordeaux.fr/~pautissi/Matrix2.pdf}
{\emph{Un lemme matriciel effectif},
Corollary 1.4}, with \(\varepsilon=1/2\).
It gives, in the height normalization
\(h_{\mathrm F}^{0}\) of (SB3),
\[
 \langle\rho_\sigma^{-2}\rangle
       \le {12\over\pi}
                 \bigl(h_{\mathrm F}^{0}(A)
                                      +{g\over2}\log(4\pi^2)\bigr)
       \le C_2(h+g).                                \tag{SB10}
\]
This theorem holds for every principally
polarized abelian variety over every number
field, with the normalized sum shown.
It has no unramifiedness, CM, fixed-degree,
or good-reduction restriction.

Set the right side of (SB9) equal to
\(M_\sigma\).
Concavity of logarithm gives
\[
 \begin{split}
 \langle\log\max(1,\rho_\sigma^{-1})\rangle
   &\le {1\over2}
                  \langle\log(1+\rho_\sigma^{-2})\rangle\\
   &\le {1\over2}\log(1+C_2(h+g)).
 \end{split}
\]
Together with (SB5) and
\(\log(1+C_2(h+g))\le C_3\log(2+g+h)\),
this proves (SB8).
Equation (SB6) proves the assertion about
all translated sections.

For the analytic content of (SB9), its
explicit proof writes \(A_\sigma\) as
\(\mathbb C^g/(\mathbb Z^g+\tau\mathbb Z^g)\),
puts \(Y=\operatorname{Im}\tau\), and uses
\[
 \|\theta\|_2=(\det(2Y))^{-1/4},\qquad
 {\|\theta\|_\infty\over\|\theta\|_2}
       \le2^{g/4}
                  \bigl(\theta(iY,0)\theta(iY^{-1},0)\bigr)^{1/2}.
\]
Poisson summation and the Gaussian
lattice-counting estimates in that proof
give (SB9).
Thus the logarithmic loss in (SB8) is a
specified theta estimate, rather than a
supremum bound depending on an unspecified
compact subset of moduli.
\end{proof}

\subsection{Integral jets and the determinant inequality}

\begin{srlemma}\label{lem:sr-slope-first-jet-integral}
Let \(x\in B\), and let \(s_1,\ldots,s_r\)
be integral sections of \(\mathcal L^n\).
There are \(L\)-linear forms \(J_i\) on
\(L^r\), \(i\ge0\), with the following
properties.
For
\[
                 F_i=\bigcap_{j<i}\ker J_j
\]
one has
\[
       J_i(F_i\cap\mathcal O_L^r)\subseteq\mathcal O_L. \tag{SB11}
\]
At every embedding, if \(f_j\) are the theta
functions of \(s_j\), and \(v_\sigma\)
is the analytic vector corresponding to
\(x\), then \(F_i\otimes_\sigma\mathbb C\)
is the space of combinations whose first
\(i\) derivatives at zero along \(v_\sigma\)
vanish.
On this space,
\[
       \sigma J_i(c)
          ={d^i\over dt^i}
                \left.\sum_j c_j f_j(tv_\sigma)\right|_{t=0}. \tag{SB12}
\]
There is no division by \(i!\) in (SB11).
\end{srlemma}

\begin{proof}
The integral Lie vector \(x\) gives an
invariant derivation on the smooth group.
Choose a local frame for \(\mathcal L\)
near zero on the generic fibre, and use the
rigidification to normalize its value at
zero.
Iterating this derivation on the local
coefficients of \(s_j\), and evaluating at
zero, defines \(J_i\).

Under a change of local frame by a unit
\(u\), the \(i\)-th derivative changes by
\(u(0)\) times the old \(i\)-th derivative
plus a linear combination of its lower
derivatives.
Thus on \(F_i\) only the value at zero of
the frame change remains.
The rigidification cancels this value.
At each finite place use a frame which is
a generator over the local ring at zero.
Its coefficients and the invariant
derivation are integral, so its first
nonvanishing derivative is integral.
The transition calculation just made
identifies it with \(J_i\).
Intersecting these valuation rings proves
(SB11).

Analytically, a local algebraic frame and an
Appell--Humbert frame differ by an invertible
holomorphic function with value one at zero
after the same rigidification.
Leibniz' rule makes their derivative
matrices triangular with diagonal one.
This proves both the kernel assertion and
(SB12).
It also shows why the individual \(J_i\)
need not be asserted integral on all of
\(\mathcal O_L^r\).
\end{proof}

\begin{srlemma}\label{lem:sr-slope-integral-flag-det}
Suppose \(J_0,\ldots,J_k\) satisfy (SB11),
and each successive form is nonzero on
\(F_i\).
Suppose positive numbers \(A_{i,\sigma}\)
satisfy
\[
 |(\sigma J_i)(c)|
      \le A_{i,\sigma}\left(\sum_j|c_j|^2\right)^{1/2}
 \quad(c\in F_i\otimes_\sigma\mathbb C).
\]
Then
\[
                   0\le\sum_{i=0}^k
                                   \langle\log A_{i,\sigma}\rangle. \tag{SB13}
\]
\end{srlemma}

\begin{proof}
The successive nonzero restrictions imply
that the \(k+1\) rows \(J_i\) are linearly
independent.
Choose a nonzero maximal minor \(\Delta\).
We first show that every such minor is
integral.
Over a finite-place DVR \(R\), the chain
\[
 R^r\supseteq F_1\cap R^r\supseteq\cdots
                      \supseteq F_{k+1}\cap R^r
\]
is a chain of saturated direct summands,
each step dropping rank by one.
Choose an adapted \(R\)-basis.
In this basis the first \(k+1\) columns
of the jet matrix are triangular and all
the other columns are zero.
Its diagonal entries are integral by
(SB11).
Its top exterior power therefore has
integral coordinates.
Changing back by a matrix in
\(\operatorname{GL}_r(R)\) preserves this
assertion, proving that \(\Delta\) is
integral at this place.
Hence \(\Delta\) is an algebraic integer.

At an archimedean place, Gram--Schmidt on
the row functionals shows that the norm of
their exterior product is at most
\(\prod_iA_{i,\sigma}\).
Indeed the \(i\)-th orthogonal row is the
restriction of \(J_i\) to the common
kernel of the preceding rows, so has norm
at most \(A_{i,\sigma}\).
Every maximal minor consequently satisfies
\[
                     |\sigma\Delta|
                            \le\prod_i A_{i,\sigma}.
\]
The product formula for the nonzero
algebraic integer \(\Delta\) gives
\(\langle\log|\sigma\Delta|\rangle\ge0\).
Taking logarithms proves (SB13).
\end{proof}

\subsection{Every jet order occurs}

\begin{srlemma}\label{lem:sr-slope-directional-jet-jumps}
Let \(\theta\ne0\) be a principal theta
function, and \(v\ne0\).
For every \(k\ge1\) and \(0\le i\le k\)
there is a theta function of level \(3k\)
which has order exactly \(i\) on the line
\(t\mapsto tv\) at zero.
\end{srlemma}

\begin{proof}
One can use the very ampleness of \(L_A^3\):
there are sections \(u_0,u_1\) of \(L_A^3\)
with \(u_0(0)\ne0\), \(u_1(0)=0\), and
\(\partial_vu_1(0)\ne0\).
Then \(u_0^{k-i}u_1^i\) has the required
order.

Here is the theta calculation producing
this precise separation, as in the source.
The functions
\[
 u_{p,q}(z)=
      \theta(z+p)\theta(z+q)\theta(z-p-q)
\]
are of level three.
There are \(p,q\) with \(u_{p,q}(0)\ne0\),
since the nonzero locus of a nonzero entire
theta function is open and dense.
If these functions failed to separate
\(v\) from the value at zero, their
logarithmic derivatives would all have
one common value at zero.
Writing
\(f=\partial_v\theta/\theta\), this says
\[
                 f(p)+f(q)+f(-p-q)=\gamma
\]
where the three denominators are nonzero.
Mixed differentiation in \(p,q\) shows
that all second derivatives of the
meromorphic function \(f\) vanish.
Thus \(f(z)\) is affine complex linear.
The theta law, however, gives
\[
                  f(z+\lambda)-f(z)
                                  =\pi H(v,\lambda).
\]
The lattice spans the tangent space over
\(\mathbb R\).
The complex-linear part on the left would
therefore equal the conjugate-linear form
\(\lambda\mapsto\pi H(v,\lambda)\) on that
whole real vector space.
Evaluation at \(v\) and \(iv\) forces
\(H(v,v)=0\), contrary to positivity.
Thus two such functions have different
ratios of derivative to value.
A linear combination gives \(u_1\) as
above, while one of them gives \(u_0\).
\end{proof}

\subsection{The directional line bound}

\begin{srproposition}\label{prop:sr-slope-integral-line-bound}
For every nonzero \(x\in B\), one has
\[
 -\left\langle\log\|x\|_\sigma\right\rangle
                    \le C_4\log(2+g+h),              \tag{SB14}
\]
where \(C_4\) is absolute.
The statement holds after every finite
extension, with the same constant.
\end{srproposition}

\begin{proof}
The case in question has \(g\ge1\).
Work first over the extension supplied by
Lemma~\ref{lem:sr-slope-joint-theta-witness};
then make the further extension for the
integer \(m\) specified below.
Let
\[
 \ell_\sigma=\|x\|_\sigma>0,\qquad n=m^2,\qquad r=m^{2g}.
\]
The field extensions preserve the normalized
degree, \(g,h\), and the normalized mean
of \(\log\|x\|\).

Choose \(m\) divisible by three, put
\(k=m^2/3\), and take the integral spanning
family \(s_P\) of (SB6).
For any complex coefficients \(c_P\), its
theta functions satisfy
\[
 \left|\sum_P c_P f_P(tv_\sigma)\right|
   \le \sqrt r\,M_\sigma
          \left(\sum_P|c_P|^2\right)^{1/2}
                    \exp\{n|t|^2\ell_\sigma^2/2\}.  \tag{SB15}
\]
Indeed each weighted norm is at most
\(M_\sigma\), Cauchy--Schwarz bounds the
sum of coefficients by their Euclidean
norm times \(\sqrt r\), and
\(\ell_\sigma^2=\pi H(v_\sigma,v_\sigma)\).

For an entire \(F\) with
\(|F(t)|\le M\exp(n|t|^2\ell^2/2)\),
Cauchy's formula on the circle
\(R=\sqrt{i/(n\ell^2)}\), \(i>0\), gives
\[
 |F^{(i)}(0)|\le i!\,M e^{i/2}
                                   (n\ell^2/i)^{i/2}
                    \le(\sqrt{enk}\,\ell)^iM
                         \qquad(1\le i\le k).
\]
The last step follows from
\((i!)^2\le i^{2i}\le(ik)^i\).
For \(i=0\), the same bound is immediate
by evaluation at zero.
Thus (SB15) gives the constants
\[
 A_{i,\sigma}
       =(\sqrt{enk}\,\ell_\sigma)^i
                                      \sqrt r\,M_\sigma. \tag{SB16}
\]
By Lemma~\ref{lem:sr-slope-first-jet-integral}
these bound the first jets on each step
of the integral vanishing filtration.
Since \(n=3k\), the spanning property and
Lemma~\ref{lem:sr-slope-directional-jet-jumps}
show that every step \(0\le i\le k\)
is nonzero.
Apply Lemma~\ref{lem:sr-slope-integral-flag-det}.
Using \(\sum_{i=0}^k i=k(k+1)/2\) gives
\[
 -\langle\log\ell_\sigma\rangle
     \le {1\over2}\log(enk)
           +{2\over k}\bigl(g\log m
                                 +\langle\log M_\sigma\rangle\bigr).
                                                               \tag{SB17}
\]
This is the explicit integral determinant
inequality; in particular, the growing
number \(m^{2g}\) of theta sections occurs
only through \(g\log m\).

Put
\[
 Y=2+g+h,\qquad
 K=\lceil2+g^2+h\rceil,\qquad m=3K,\qquad k=3K^2.
\]
Then \(Y\ge3\), \(m\le3Y^2\), and
\(k\ge g^2+h+1\).
Consequently
\[
 \log m\le3\log Y,\qquad
 {g\over k}\le1,\qquad {h\over k}\le1,\qquad
 {1\over2}\log(enk)
        ={1\over2}\log(3e)+\log k\le C_5\log Y.
\]
Insert (SB8) in (SB17).
Each remaining term is at most an absolute
constant times \(\log Y\).
This proves (SB14), with all constants
chosen before \(A,L,g\) and before any
extension.

Finally the vector \(x\) here is the base
change of the original vector.
For every finite \(L'/L\),
\[
 {1\over[L':\mathbb Q]}
     \sum_{\tau:L'\hookrightarrow\mathbb C}
            \log\|x\|_\tau
       ={1\over[L:\mathbb Q]}
     \sum_{\sigma:L\hookrightarrow\mathbb C}
            \log\|x\|_\sigma,
\]
because each \(\sigma\) has \([L':L]\)
extensions.
Thus the argument proves (SB14) over the
original field as well.
\end{proof}

\subsection{The absolute line theorem and maximal slope}

The precise absolute Siegel input is:
for every Hermitian bundle \(\overline V\)
of rank \(d>0\), and every \(\varepsilon>0\),
there are a finite extension and a line
subbundle \(M\) of its base change with
\[
 \widehat{\deg}_{\mathrm n}\overline M
       \ge\widehat\mu(\overline V)
                              -\tfrac12\log d-\varepsilon. \tag{SB18}
\]
The line and the extension may depend on
\(\overline V,\varepsilon\).
This is a consequence of Zhang's absolute
successive-minima theorem, in the form of
Bost--Chen,
\href{https://webusers.imj-prg.fr/~huayi.chen/Recherche/bost_chen.pdf}
{\emph{Concerning the semistability of
tensor products in Arakelov geometry},
Theorem 2.5 and Corollary 2.7}.
Indeed the corollary gives spanning lines
whose sum of degrees is at least
\(\widehat{\deg}_{\mathrm n}\overline V
-\tfrac d2\sum_{j=2}^d j^{-1}-\varepsilon\);
one line has at least the average degree,
and \(\sum_{j=2}^d j^{-1}\le\log d\).
Every line over \(\overline{\mathbb Q}\)
appearing here descends to a finite
extension. This is the scope of the
absolute theorem; no assertion about a
short vector over the original number
field is used.

\begin{proof}[Proof of Theorem~\ref{thm:sr-general-principal-slope}]
If \(g=0\), the Lie bundle has no nonzero
submodule and the inequality is immediate.
Assume \(g\ge1\).
Let \(F\subseteq B\) be a nonzero saturated
submodule of rank \(d\), with its induced
metric.
Apply (SB18) to \(F\), with
\(\varepsilon=1\).
After a further finite extension,
principalize the underlying ideal of the
resulting saturated line \(M\).
An integral generator \(x\ne0\) of that
line then satisfies
\[
 -\langle\log\|x\|_\sigma\rangle
       =\widehat{\deg}_{\mathrm n}\overline M
       \ge\widehat\mu(\overline F)
                                     -\tfrac12\log d-1. \tag{SB19}
\]
The equality uses an actual generator,
so has no extra finite-place index.
Normalized degree is unchanged by
principalization and base change.

The base change of \(B\) is the actual Lie
bundle, and its metrics are isometric,
by Lemma~\ref{lem:sr-slope-model-metric}.
Proposition~\ref{prop:sr-slope-integral-line-bound}
therefore applies to this \(x\).
Since \(d\le g\), (SB19) gives
\[
 \widehat\mu(\overline F)
       \le C_4\log(2+g+h)+\tfrac12\log g+1
       \le(C_4+2)\log(2+g+h).
\]
Taking the supremum over \(F\) proves (SB1).
Nonsaturated submodules have no larger
degree than their saturation.
The endomorphism host was never used in
choosing any constant.
\end{proof}

\subsection{Exact relation to the frozen producer chain}

The final determinant and line-to-slope
steps are the mathematical content of
\[
 \begin{gathered}
 \srcid{ExistsIntegralJetBasis.jets},\qquad
 \srcid{IntegralFlagEmbeddingCriterion.flag\_slope},\\
 \srcid{ThetaDerivativeBound.lefschetz\_direction},\qquad
 \srcid{LineOfKeyDegree.core\_applied},\\
 \srcid{LineOfKeyDegree.line\_of\_key},\qquad
 \srcid{SlopeBoundOfLargeDegreeLine.slopeBound}.
 \end{gathered}
\]
The first five modules have prefix
\srcid{Development.AutomorphicSpectrum.Proofs.};
their respective subdirectories are
\srcid{GeometryProofs}, \srcid{AdelicProofs},
\srcid{MetricProofs}, \srcid{GeometryProofs},
\srcid{GeometryProofs}.
The last module is in
\srcid{GeometryProofs}.
The frozen blobs for the last two are
\srcid{666e2dd3ab1ce84b3448899b33f83d08030c4ad3}
and
\srcid{89a20c55b564e35e40ef9c2a02178c32b8104a8c}.

The source obtains the theta mean input by
\[
 \srcid{KeyBadReductionHolds.key\_bad}
 \longrightarrow
 \srcid{KeyDegreeOfBadReduction.keyDeg\_of\_keyBad}.
\]
The latter has frozen blob
\srcid{9ac84964f5bca6b0f36e6538f59ba6a393efb2ea}.
Its key scalar satisfies, for an integral
Lie basis \(e\),
\[
 b_\sigma^8=
 { (2\pi)^{2g}
             |\det\operatorname{Gram}_{B,\sigma}(e)|^2
        \over|\sigma\kappa|}.                         \tag{SB20}
\]
Together with
\(\langle\log|\det\operatorname{Gram}(e)|\rangle
=-2\widehat{\deg}_{\mathrm n}\overline B\),
the integrality of \(\kappa\ne0\) gives
precisely the upper bound in (SB5).
The bad-reduction producer runs through
\srcid{WeilBound.CoreStatements.HodgeTatePoints.LieMultComp},
whose frozen blob is
\srcid{3f7d4370da34864229be0897f207d911de688a71}.
It squares the packet, places it in a
unitary moduli problem, and invokes
\srcid{KftwoModel.key\_unit}.
The latter uses
\srcid{ThetaKeyScalarIsUnitOnIntegralModels}:
normal charts carry the key function and
its inverse; good-reduction trait values
are units; degeneration arcs have bounded
theta coefficient and nonzero limiting
Lie determinant; the no-pole criterion and
normality then extend both functions.
This is a genuine integral-model argument,
not an inference from complex theta
identities alone.

Likewise the archived
\srcid{TorsionValueIntegral} proves
integrality of its specific normalized
torsion scalar through Mumford
uniformization and a weight-space valuation
bound; the latter uses shortest
representatives of the quadratic form.
That universal assertion must not be
replaced by the claim that every torsion
point extends to the connected model.

The present proof supplies instead the
joint witness of
Lemma~\ref{lem:sr-slope-joint-theta-witness}
from the cited classical model theorem,
with its finite-place compatibility
proved above.
Equations (SB5)--(SB6) are exactly the
theta inputs actually used in
(SB15)--(SB17).
The explicit matrix and Gromov estimates
supply (SB8).
Consequently no uninstantiated
\(\mathrm{KeyBadSt}\),
\(\mathrm{KeyDegSt}\), or
theta-supremum hypothesis remains in
Theorem~\ref{thm:sr-general-principal-slope}.
We have not thereby supplied a separate
prose proof of the universal key-rule and
torsion-value declarations for arbitrary
inputs: their stronger quantifiers are
distinct from the joint witness sufficient
for this theorem.
The published inputs retained here are
the precise Moret--Bailly model theorem,
the two specified principal-theta
estimates, the Neron-model comparison, and
Zhang's absolute line theorem.
